\documentclass[11pt,a4paper]{article}
\usepackage[utf8]{inputenc}
\usepackage[T1]{fontenc}
\usepackage{amsmath,amssymb,amsthm}
\usepackage[margin=2.5cm]{geometry}
\usepackage{graphicx}
\usepackage{enumitem}
\usepackage{xcolor}
\usepackage[most]{tcolorbox}
\usepackage{titlesec}
\usepackage{fancyhdr}
\definecolor{noteblue}{HTML}{185FA5}
\definecolor{notebluebg}{HTML}{E6F1FB}
\definecolor{noteteal}{HTML}{0F6E56}
\definecolor{notetealbg}{HTML}{E1F5EE}
\definecolor{noteamber}{HTML}{854F0B}
\definecolor{noteamberbg}{HTML}{FAEEDA}
\definecolor{notepurple}{HTML}{534AB7}
\definecolor{notepurplebg}{HTML}{EEEDFE}
\definecolor{notered}{HTML}{A32D2D}
\definecolor{noteredbg}{HTML}{FCEBEB}
\definecolor{notegray}{HTML}{444441}
\definecolor{notegraybg}{HTML}{F1EFE8}
\titleformat{\section}{\Large\bfseries\color{noteblue}}{\thesection}{1em}{}[{\color{noteblue}\titlerule[1.2pt]}]
\titleformat{\subsection}{\large\bfseries\color{noteblue}}{\thesubsection}{1em}{}
\newtcbtheorem[number within=section]{definizione}{Definizione}{enhanced,breakable,colback=notebluebg,colframe=noteblue,colbacktitle=notebluebg,coltitle=noteblue,fonttitle=\bfseries,boxrule=0pt,leftrule=3pt,titlerule=0pt,arc=0pt}{def}
\newtcbtheorem[number within=section]{teorema}{Teorema}{enhanced,breakable,colback=notetealbg,colframe=noteteal,colbacktitle=notetealbg,coltitle=noteteal,fonttitle=\bfseries,boxrule=0pt,leftrule=3pt,titlerule=0pt,arc=0pt}{teo}
\newtcbtheorem[number within=section]{esempio}{Esempio}{enhanced,breakable,colback=noteamberbg,colframe=noteamber,colbacktitle=noteamberbg,coltitle=noteamber,fonttitle=\bfseries,boxrule=0pt,leftrule=3pt,titlerule=0pt,arc=0pt}{ex}
\newtcolorbox{intuizione}{enhanced,breakable,colback=notepurplebg,colframe=notepurple,colbacktitle=notepurplebg,coltitle=notepurple,fonttitle=\bfseries,boxrule=0pt,leftrule=3pt,titlerule=0pt,arc=0pt,title=Intuizione}
\newtcolorbox{attenzione}{enhanced,breakable,colback=noteredbg,colframe=notered,colbacktitle=noteredbg,coltitle=notered,fonttitle=\bfseries,boxrule=0pt,leftrule=3pt,titlerule=0pt,arc=0pt,title=Attenzione}
\newtcolorbox{sintesi}{enhanced,breakable,colback=notegraybg,colframe=notegraybg,colbacktitle=notegraybg,coltitle=notegray,fonttitle=\bfseries,boxrule=0pt,titlerule=0pt,arc=2pt,title=In sintesi}
\theoremstyle{plain}
\newtheorem{lemma}{Lemma}[section]
\newtheorem{corollario}[lemma]{Corollario}
\pagestyle{fancy}
\fancyhf{}
\fancyhead[L]{\small\itshape Model Order Reduction}
\fancyhead[R]{\small\itshape 12-06-2026}
\fancyfoot[C]{\small--- \thepage\ ---}
\renewcommand{\headrulewidth}{0.4pt}
\setlength{\headheight}{14pt}

\title{Model Order Reduction}
\date{12-06-2026}
\author{}
\begin{document}
\maketitle
\section{La Proper Orthogonal Decomposition}
Paragrafo di prosa introduttiva con matematica inline $u(\mu)\in\mathbb{R}^N$ e un'equazione:
\begin{equation}
S = [\,u(\mu_1)\,|\,\dots\,|\,u(\mu_{n_s})\,]
\end{equation}
\begin{definizione}{Matrice degli snapshot}{snapshot}
Le colonne di $S$ sono le soluzioni del problema completo.
\end{definizione}
Altra prosa tra i box, come richiesto dalla regola 14.
\begin{intuizione}
Le fotografie-base con cui ricostruire tutte le altre.
\end{intuizione}
\subsection{Quanti modi tenere?}
Prosa con casi:
\[ \sigma_i \approx \begin{cases} 1 & i \le N \\ 0 & i > N \end{cases} \]
\begin{esempio}{Equazione del calore}{calore}
Con $N=5$ modi l'errore scende sotto $10^{-4}$.
\end{esempio}
Prosa di raccordo.
\begin{teorema}{Ottimalit\`a della POD}{pod-opt}
La base POD minimizza l'errore quadratico di proiezione.
\end{teorema}
\begin{lemma}
Un risultato secondario senza box.
\end{lemma}
\begin{corollario}
E il suo corollario.
\end{corollario}
\begin{attenzione}
Nei problemi advection-dominated la POD da sola non basta.
\end{attenzione}
Riferimenti incrociati: per la Definizione~\ref{def:snapshot} e il Teorema~\ref{teo:pod-opt}, vedi l'Esempio~\ref{ex:calore}.
\begin{sintesi}
\begin{itemize}
\item Gli snapshot raccolgono soluzioni note.
\item Il numero di modi si sceglie dal decadimento dei valori singolari.
\end{itemize}
\end{sintesi}
\end{document}
